\chapter{Le Web, et ce qui est déclaré sans être éprouvé}

Deux sujets dans ce chapitre, et ils vont ensemble : une cible qui fonctionne et
qu'on connaît mal, et six cibles dont il faut dire honnêtement ce qu'on en sait.

\section{Le Web}

\begin{nkcode}{Applications/NkDames/NkDames.jenga -- reel}
with filter("system:Web"):
    consoleapp()
    usetoolchain("emscripten")
    defines(["NKENTSEU_ENABLE_EMSCRIPTEN"])
    emscriptencanvasid("canvas")
    emscripteninitialmemory(32)
    emscriptenextraflags([
        "-s", "ASYNCIFY",
        "-s", "ALLOW_MEMORY_GROWTH=1",
        "-s", "FULL_ES3=1",
        "-s", "USE_WEBGL2=1",
        "-s", "NO_EXIT_RUNTIME=1",
    ])
\end{nkcode}

\begin{nkretenir}
\textbf{\texttt{consoleapp()} sur le Web, alors que le bureau dit
\texttt{windowedapp()}.} Ce n'est pas une erreur : la notion de fenêtre native
n'existe pas dans un navigateur, la page \emph{est} la fenêtre.

C'est un bon exemple de ce que les filtres servent à faire --- redéfinir même le
genre du projet selon la cible.
\end{nkretenir}

\begin{center}
\begin{tabular}{@{}ll@{}}
\toprule
\textbf{Drapeau} & \textbf{Ce qu'il change} \\
\midrule
\texttt{ASYNCIFY} & permet à la boucle de rendre la main \\
\texttt{ALLOW\_MEMORY\_GROWTH} & la mémoire n'est plus figée au démarrage \\
\texttt{FULL\_ES3} + \texttt{USE\_WEBGL2} & la traduction OpenGL ES vers WebGL2 \\
\texttt{NO\_EXIT\_RUNTIME} & le programme ne se termine pas à la fin du \texttt{main} \\
\bottomrule
\end{tabular}
\end{center}

\begin{nkretenir}[\texttt{ASYNCIFY} n'est pas une option de confort]
Un navigateur exécute votre page sur \emph{son} fil principal --- celui qui
dessine et qui répond aux clics. Une boucle infinie n'y rend jamais la main :
l'onglet gèle, et le navigateur finit par proposer de le tuer.

\texttt{ASYNCIFY} rend possible de suspendre et reprendre le programme, ce qui
permet à la boucle de céder. \textbf{Sans lui, votre programme est correct et
votre onglet est mort.}

\alerte{} Et le drapeau ne suffit pas : encore faut-il que la boucle \emph{cède
réellement}. Un mécanisme de temporisation qui existe et qu'on n'appelle pas
laisse le même symptôme.
\end{nkretenir}

\begin{nkpiege}
\textbf{Un \texttt{.wasm} ouvert en double-cliquant le \texttt{.html} ne se
charge pas.} Les navigateurs refusent le protocole \texttt{file://} pour ce genre
de ressource.

Il faut \emph{servir} le dossier par HTTP --- \texttt{python -m http.server}
suffit. Le message d'erreur, lui, parle de \emph{CORS} ou de \emph{fetch}, et
n'aide pas à trouver la cause.
\end{nkpiege}

\begin{nkretenir}
\texttt{emscripteninitialmemory(32)} donne 32 Mo au départ. Avec
\texttt{ALLOW\_MEMORY\_GROWTH}, ce n'est qu'un point de départ --- mais un point
de départ trop petit fait payer plusieurs réallocations pendant le chargement,
qui est justement le moment où l'utilisateur regarde.
\end{nkretenir}

\section{Ce qui est déclaré, et ce qui est éprouvé}

C'est la section la plus honnête du livre, et la plus utile avant une promesse.

\begin{center}
\begin{tabular}{@{}lll@{}}
\toprule
\textbf{Cible} & \textbf{État} & \textbf{Ce qui le fonde} \\
\midrule
Windows & \textbf{éprouvée} & artefacts construits et lancés ici \\
Linux & \textbf{éprouvée} & construits par WSL2, binaires lancés \\
Web & \textbf{éprouvée} & \texttt{.wasm} produit, page servie \\
Android & \textbf{éprouvée} & APK installé sur un téléphone réel \\
HarmonyOS & construite & \texttt{.hap} produit, non installé ici \\
macOS & déclarée & exige macOS + Xcode --- CI \\
iOS & déclarée & idem, et le descripteur le dit \\
PS4, PS5, Xbox, Switch & déclarées & SDK sous accord, non testés ici \\
FreeBSD, OpenBSD & déclarées & aucune mesure \\
\bottomrule
\end{tabular}
\end{center}

\begin{nkretenir}
\textbf{Trois états, et il faut les trois mots.}

\emph{Déclarée} : des lignes existent dans un descripteur et une chaîne d'outils
est prévue. Rien ne dit qu'elles fonctionnent.

\emph{Construite} : un artefact sort. Rien ne dit qu'il se lance.

\emph{Éprouvée} : l'artefact a tourné sur la cible.

\textbf{Annoncer « douze plateformes » sans cette distinction est une promesse
qu'on ne peut pas tenir.} Et l'écart n'est pas théorique : le filtre iOS de ce
dépôt \textbf{n'existait pas} alors que l'en-tête du fichier annonçait iOS depuis
des mois.
\end{nkretenir}

\begin{nkpiege}
\textbf{Un commentaire qui nomme une plateforme n'est pas un support.} Le cas est
mesuré : l'en-tête disait « Mobile : Android, iOS, HarmonyOS », et aucun filtre
iOS n'existait dans le fichier.

Rien ne le signalait : \textbf{la construction bureau ne dit rien d'un filtre
manquant.} On croit couvrir une plateforme parce qu'une phrase la nomme.

Le contrôle est immédiat : \texttt{grep "system:iOS"} dans le descripteur. S'il
n'y a pas de filtre, il n'y a pas de cible.
\end{nkpiege}

\section{Les consoles}

\begin{nkretenir}
Jenga déclare PS4, PS5, Xbox One, Xbox Series et Switch, avec leurs fonctions :
\texttt{gdkpath}, \texttt{xboxmode}, \texttt{xboxpackagename},
\texttt{xboxsigningmode}\dots

\textbf{Leurs SDK ne sont ni publics ni redistribuables.} Le code qui les gère
existe dans \texttt{Jenga/Core/Builders/}, il n'a pas pu être exercé ici, et ce
livre ne prétendra pas le contraire.

Ce que vous pouvez en faire aujourd'hui : lire les descripteurs comme un modèle
de ce qu'une cible sous accord demande --- une identité de paquet, un mode de
signature, un chemin de SDK.
\end{nkretenir}

\section{Choisir sa cible en ligne de commande}

\begin{nkcode}{}
jenga build --platform windows
jenga build --platform linux --linux-backend=wayland
jenga build --platform android --config Release
jenga build --platform web
\end{nkcode}

\begin{nkretenir}
\textbf{Sans \texttt{-{}-platform}, Jenga construit pour la machine courante.}
C'est le défaut le plus sûr, et c'est aussi pourquoi on oublie qu'il existe : on
peut croire tester la portabilité en ne construisant jamais que pour soi.
\end{nkretenir}

\begin{nkbilan}
Le Web se décrit comme une console --- la page est la fenêtre --- et exige
\texttt{ASYNCIFY}, sans quoi la boucle gèle l'onglet ; son artefact doit être
servi par HTTP, jamais ouvert en \texttt{file://}. Les douze cibles déclarées ne
sont pas au même stade : \emph{déclarée}, \emph{construite} et \emph{éprouvée}
sont trois états distincts, et les confondre est une promesse qu'on ne tient pas.
Un commentaire qui nomme une plateforme n'est pas un support --- mesuré sur ce
dépôt, où iOS était annoncé sans qu'aucun filtre n'existe.
\end{nkbilan}

\begin{nkquestions}
\begin{enumerate}
\item Pourquoi le Web est-il déclaré \texttt{consoleapp()} ?
\item Que se passe-t-il sans \texttt{ASYNCIFY} ?
\item Pourquoi un \texttt{.wasm} ne se charge-t-il pas en \texttt{file://} ?
\item Distinguez « déclarée », « construite » et « éprouvée ».
\item Comment vérifier en une commande qu'une plateforme est réellement décrite ?
\end{enumerate}
\end{nkquestions}

\begin{nkpratique}
Construisez votre projet pour le Web, puis servez-le et ouvrez-le.

\begin{enumerate}
\item \texttt{jenga build -{}-platform web} ;
\item trouvez les fichiers produits --- combien, et lesquels ? ;
\item ouvrez le \texttt{.html} en double-cliquant, et notez le message ;
\item servez le dossier par HTTP, ouvrez, et notez la différence.
\end{enumerate}

Le point 3 est le plus formateur : \textbf{gardez le message}, vous le
reverrez.
\end{nkpratique}

\begin{nkdemo}
Prouvez qu'un commentaire n'est pas un support.

Prenez un descripteur du dépôt et dressez deux listes :
\begin{enumerate}
\item les plateformes que son \textbf{en-tête} annonce ;
\item les plateformes pour lesquelles un \texttt{with filter("system:\dots")}
      existe réellement.
\end{enumerate}

\textbf{Ce que la démonstration établit} : les deux listes peuvent différer, et
rien dans la construction ne le signale.

Refaites l'exercice sur \emph{votre} projet dès qu'il annonce une plateforme.
C'est un contrôle d'une minute qui empêche une promesse fausse.
\end{nkdemo}

\begin{nklogique}
On vous demande : « le projet supporte-t-il macOS ? »

\begin{enumerate}
\item Quelles trois réponses différentes pourriez-vous donner, selon l'état réel ?
\item Pour chacune, quelle preuve produiriez-vous ?
\item L'une de ces preuves ne peut pas être produite sur votre machine.
      Laquelle, et que répondez-vous alors --- sans mentir ni renoncer ?
\end{enumerate}
\end{nklogique}
